% Source: https://artofproblemsolving.com/community/c3h2567636_number_theory
% Stable identity: aops-topic-2567636-q01
\textbf{Existence} (strong induction on $n$). For $n=1$ take $1=1\cdot3^0$. Let $n\ge2$ and let $r\in\{-1,0,1\}$ with $r\equiv n\pmod3$. Then $m=\dfrac{n-r}{3}$ is an integer with $1\le m<n$ (as $n\ge2$). By the induction hypothesis $m=\sum_{j=0}^{t}e_j3^j$ with $e_j\in\{-1,0,1\}$ and $e_t\ne0$, so
\[n=3m+r=\sum_{j=0}^{t}e_j3^{j+1}+r\cdot3^0\]
is a representation of the required form, with leading digit $e_t\ne0$.

\textbf{Uniqueness.} Suppose $\sum c_j3^j=\sum d_j3^j$ are two different representations (pad the shorter one with zeros). Let $k$ be the smallest index with $c_k\ne d_k$. Then $\sum_{j\ge k}(c_j-d_j)3^j=0$; dividing by $3^k$ shows $3\mid c_k-d_k$. But $c_k-d_k\in\{-2,-1,1,2\}$, a contradiction. So the digits agree, and since both leading digits are nonzero, the lengths agree too. $\blacksquare$

\textbf{The student's argument.} Replacing $2\cdot3^j$ by $3^{j+1}-3^j$ adds $1$ to the digit in position $j+1$, which can then become $2$ or $3$. For example $5=1\cdot3+2$ becomes $1\cdot3+3-1=2\cdot3-1$, which still has a digit $2$; one more step gives $5=9-3-1$. So the replacement must be repeated with carrying, and one has to show that the process stops. The argument also says nothing about uniqueness.
